Automated theorem proving for first-order logic can be studied
through sequent calculi, where a proof is built by repeatedly
decomposing a goal sequent into simpler subgoals. This report
investigates the backward proof-search algorithm for the pedagogical
LK$'$ calculus, starting with a na\"ive backtracking prover that,
while simple to implement and closely following the calculus, has a
search space that grows quickly under rule choices, branching
connectives, and quantified formulae. A Rust implementation was
developed to evaluate two refinements: iterative deepening, which
repeatedly runs depth-bounded search to find shallow proofs before
deeper distractions, and a six-class rule priority schedule that
attempts closing and low-branching rules before more expansive
quantifier steps. On a 737-problem provable subset of the TPTP
library~\cite{Sut17}, the baseline closes 152 problems within a
one-second budget, iterative deepening closes 199, and the combined
priority-ID prover closes 209---demonstrating that small control
decisions matter significantly in proof search. While the logical
rules define what is sound, the search strategy determines whether a
proof is found within practical limits; iterative deepening produces
most of the improvement by avoiding premature commitment to deep
branches, while priority scheduling gives consistent but smaller
benefits by delaying costly rule applications.

\keywords{Sequent calculus \and LK$'$ \and Backward proof search \and
Iterative deepening \and First-order logic \and TPTP.}
